\chapter{Related Work}

\section{Dark Pattern Taxonomies}

Dark patterns, also referred to as deceptive user interfaces, have been widely studied as recurring interface strategies that steer, coerce, or deceive users into decisions that they may not otherwise have made. Early work in human-computer interaction frames dark patterns not merely as examples of poor usability, but as ethically problematic design practices that redirect user attention, effort, or interpretation toward outcomes preferred by service providers \cite{gray2018darkpatterns}. This perspective is important because it shifts the discussion from accidental interface inconvenience to systematic manipulation. A confusing button, a hard-to-find cancellation path, or a visually dominant purchase option may appear as an isolated design choice when observed individually. However, when such patterns recur across platforms and are organized into stable categories, they become measurable design phenomena.

Prior taxonomies have identified a range of dark pattern types, including obstruction, interface interference, sneaking, forced action, nagging, scarcity, social proof, hidden costs, and difficult cancellation \cite{gray2018darkpatterns,mathur2019darkpatterns,mathur2021whatmakesdark}. Obstruction refers to designs that make a user’s intended action unnecessarily difficult, such as requiring many steps to cancel a subscription or reject an optional service. Interface interference changes the presentation of choices by visually emphasizing one option while hiding or weakening another. Sneaking occurs when costs, commitments, or consequences are introduced without sufficient salience, often becoming visible only near the end of a transaction. Forced action requires users to complete an unrelated or undesired action before accessing a service, while nagging repeatedly interrupts the user with prompts or requests that are difficult to dismiss. Scarcity and social proof create pressure by suggesting that a product is limited or popular, while hidden costs and difficult cancellation exploit information asymmetry and procedural friction.

These taxonomies are useful for two reasons. First, they provide a vocabulary for describing deceptive interfaces with greater precision. Instead of treating all manipulative designs as a single vague category, taxonomies distinguish between different mechanisms of influence. Second, they show that dark patterns are systematic rather than isolated design mistakes. If multiple platforms use similar mechanisms, then the problem cannot be fully explained by individual design errors. It reflects broader incentives in digital markets, where interfaces are optimized not only for usability but also for conversion, retention, monetization, and data collection. This report builds on this taxonomic tradition, but it does not attempt to reproduce every fine-grained category. Instead, it operationalizes dark pattern exposure using a seven-dimensional DPS vector that captures several recurring mechanisms: Interface Interference, Friction, Pressure, Contextual Deception, Sneaking, Nagging, and Forced Action.

\section{Dark Patterns as Choice Architecture Manipulation}

A second line of work explains dark patterns through the concept of choice architecture. Choice architecture refers to the way options are organized, displayed, sequenced, and made actionable within a decision environment. In digital interfaces, user decisions are rarely made in a neutral space. Buttons have different colors and sizes, options appear in different locations, defaults may already be selected, and some actions require more effort than others. As a result, the interface can shape decision-making even when the formal set of available options remains unchanged \cite{mathur2021whatmakesdark}.

This perspective is especially useful for understanding why dark patterns can be harmful without relying on a narrow definition of deception. A design does not need to contain a false statement in order to manipulate user choice. It may instead alter option visibility, default choices, visual emphasis, information flow, or the relative effort required to accept or reject an option. For example, an application may make the purchase button large and colorful while placing the decline option in small grey text. It may allow users to accept a membership offer with one tap but require several screens to refuse or cancel it. It may reveal fees only after the user has invested time in the transaction. These designs influence user behavior by restructuring the decision environment.

This interpretation directly motivates several DPS dimensions used later in this report. Interface Interference captures cases where visual hierarchy, layout, or salience directs attention toward a particular option. Friction captures asymmetries in effort, especially when rejecting, leaving, or cancelling is made harder than accepting. Pressure captures urgency, scarcity, countdowns, emotional prompts, and similar cues that encourage quick action. Forced Action captures situations in which users must complete an action, such as registration, payment binding, app installation, or permission granting, before they can proceed. In this sense, the DPS vector is not merely a list of labels. It is an attempt to translate the broader theory of choice architecture manipulation into observable annotation dimensions suitable for screenshot-level and interaction-level auditing.

Understanding dark patterns as choice architecture manipulation also helps avoid overclaiming about user intention or platform intent. This report does not assume that every observed design was intentionally created to deceive. Instead, it focuses on observable interface mechanisms and their association with different user profiles. The question is not whether a designer had malicious intent, but whether controlled profiles encounter systematically different forms or levels of manipulative interface exposure.

\section{Large-scale Measurement of Deceptive Interfaces}

Large-scale measurement studies have shown that dark patterns can be empirically observed, classified, and quantified across digital services. Prior work has examined deceptive patterns in online shopping websites, consent interfaces, and other web or app environments, demonstrating that manipulative designs are not limited to anecdotal examples \cite{mathur2019darkpatterns,nouwens2020gdpr,mathur2021whatmakesdark}. These studies are important because they transform dark patterns from a primarily conceptual or normative concern into a measurable empirical object. Once researchers can identify and count dark pattern instances across many interfaces, they can compare prevalence, categorize mechanisms, and discuss platform accountability with stronger evidence.

For example, large-scale studies of shopping interfaces have shown that commercial platforms frequently use designs related to urgency, scarcity, social proof, hidden costs, and obstruction \cite{mathur2019darkpatterns}. Work on consent pop-ups further demonstrates how interface details such as button placement, default settings, and the visibility of opt-out options can affect user choices \cite{nouwens2020gdpr}. Such studies support the view that interface design can systematically influence user autonomy and consent. They also provide methodological inspiration for this report: dark patterns can be annotated, counted, and compared rather than discussed only through individual examples.

However, much of this measurement work focuses on the presence or absence of dark patterns within an interface. In other words, the central question is often whether a website, app, or page contains a dark pattern. This is a necessary foundation, but it does not fully capture the experience of users in personalized and interactive systems. Mobile applications are not static documents. Users move through different screens, search for different keywords, click different items, and trigger different recommendation or monetization flows. Therefore, two users may interact with the same application but encounter different prompts, pop-ups, membership offers, payment pages, or cancellation paths.

This distinction motivates the shift from dark pattern presence to dark pattern exposure. Presence asks whether a dark pattern exists somewhere in the system. Exposure asks whether a particular user profile actually encounters it during interaction. A dark pattern that exists deep inside an application may have little immediate effect if most users never reach it. Conversely, a pattern that repeatedly appears along a common interaction path may strongly shape user experience. Existing measurement work provides the tools for identifying deceptive interfaces, but this report extends the question toward profile-conditioned exposure along interaction trajectories.

\section{Personalization and Algorithm Auditing}

Digital systems increasingly personalize the outputs that users receive. Search results, recommendations, advertisements, rankings, prices, feeds, notifications, and promotional content may vary according to user history, inferred interests, behavior, device context, or account state. Prior research on personalization has shown that different users can receive different outputs even when they interact with the same platform or issue similar queries \cite{hannak2013personalization}. This finding has motivated broader concerns about differential exposure, because personalized systems may shape what users see, what they believe is available, and what actions they are encouraged to take.

Algorithm auditing provides a methodological foundation for studying such systems from the outside. Because platform internals are often inaccessible, researchers commonly use controlled accounts, scripted behaviors, or sock-puppet profiles to observe outputs under different conditions \cite{sandvig2014auditing}. Auditing studies of recommendation systems, search engines, and content platforms have examined whether different profiles receive different rankings, recommendations, advertisements, or problematic content \cite{ribeiro2020auditing,haroon2023youtubeaudit,srba2023youtube}. These works show that profile-based comparison can reveal differences that would remain hidden if researchers only inspected a platform from a single account or a single interaction path.

This report applies the logic of algorithm auditing to deceptive interface exposure. Instead of asking whether different profiles receive different videos, search results, or advertisements, it asks whether they receive different dark pattern exposure during mobile app interaction. This is a natural extension of auditing research because modern applications often combine personalization with interface optimization. A user who frequently searches for discounts may be shown more promotional prompts. A user who repeatedly enters payment pages may encounter more membership offers, countdowns, or checkout-related pressure. A user who mostly browses without purchasing may encounter fewer payment-oriented interventions. These possibilities should be treated carefully as empirical questions rather than assumed outcomes. Nevertheless, prior auditing research motivates the hypothesis that dark pattern exposure may also be profile-dependent.

The connection between personalization and dark patterns is especially important for social network security and privacy. If manipulative interfaces are personalized, then their impact may not be evenly distributed across users. Some profiles may be more frequently exposed to pressure, forced action, or obstructive flows than others. Ignoring this possibility may obscure how certain behavioral groups are disproportionately steered toward monetization, data disclosure, or continued engagement. Therefore, a profile-aware audit does not replace traditional dark pattern detection. It complements it by asking who encounters which patterns, under what interaction conditions, and with what degree of consistency.

\section{User Modeling for Controlled Audits}

User modeling research commonly represents users through interests, behaviors, and attributes \cite{eke2019userprofiling,piao2018interests,purificato2024usermodeling}. Interests describe what users like or seek, such as product categories, topics, content genres, or keywords. Behaviors describe how users interact, including searching, browsing, clicking, liking, following, watching, staying, dismissing, or entering decision pages. Attributes describe who users are, such as demographic or relatively stable personal characteristics. In real-world personalization systems, these dimensions may be combined to infer preferences and predict future actions.

For controlled audits, however, the goal is not to reconstruct real users in full detail. The goal is to create reproducible profiles that differ along selected dimensions while minimizing uncontrolled variation. This report therefore models user profiles through interests and behaviors, while excluding demographic attributes. Interests are operationalized as probability distributions over keywords. For example, a high-payment shopping profile can be associated with higher probability for high-value product keywords, while a bargain-hunter profile can be associated with discount-related keywords. This design is conceptually related to probabilistic interest modeling, where latent preferences are represented through observable topics, words, or behavioral traces \cite{blei2003lda,piao2018interests}. The report does not claim to infer real user interests through a trained model. Rather, it uses manually specified keyword distributions to create controlled and repeatable audit conditions.

Behaviors are modeled as probability distributions over interaction actions. This abstraction is also consistent with prior work in user behavior modeling and recommendation systems, where clicks, browsing, searches, and sequential actions are treated as behavioral signals \cite{chapelle2009clickmodel,zhou2018din,li2019mind}. In this project, behavior modeling is simplified for auditability. Each controlled profile has an action distribution that determines how often it searches, browses, visits decision interfaces, or performs other interaction steps. The purpose is not to simulate the complete complexity of real mobile users, but to create distinct and interpretable interaction trajectories.

The exclusion of demographic attributes is a deliberate design choice. Demographic features are difficult to control reliably in black-box app audits, especially when platforms may infer them through indirect signals rather than explicit profile fields. They may also raise ethical concerns, particularly if the audit creates or manipulates sensitive identities. By focusing on interests and behaviors, this project studies differential exposure through controllable and reproducible profile dimensions. This choice does not imply that demographic differences are unimportant. Rather, it reflects a course-project scope and an ethical preference for avoiding sensitive attributes while still examining profile-conditioned exposure.

\section{Research Gap}

The literature reviewed above suggests three important conclusions. First, dark pattern research has developed rich taxonomies and measurement methods for identifying deceptive interfaces \cite{gray2018darkpatterns,mathur2019darkpatterns,mathur2021whatmakesdark,nouwens2020gdpr}. These works show that dark patterns are systematic, classifiable, and empirically measurable. Second, algorithm auditing research has demonstrated that digital platforms can produce different outputs for different controlled profiles, especially in personalized search, recommendation, advertising, and content-ranking environments \cite{sandvig2014auditing,hannak2013personalization,ribeiro2020auditing}. Third, user modeling research provides practical abstractions for constructing profiles through interests and behaviors, which can be adapted for controlled audit experiments \cite{eke2019userprofiling,purificato2024usermodeling}.

However, these lines of work are rarely connected directly. Existing dark pattern studies often ask whether a deceptive pattern exists in an interface, but they do not always ask whether different users are exposed to different patterns through different interaction paths. Algorithm auditing studies examine differential outputs across profiles, but they usually focus on search results, recommendations, advertisements, or content exposure rather than deceptive interface mechanisms. User modeling studies provide ways to represent users, but they are not usually used to measure profile-dependent dark pattern exposure.

This report addresses that gap by proposing a profile-aware auditing framework for mobile applications. The framework treats dark pattern exposure as a profile-conditioned outcome rather than a static property of an app. It constructs controlled profiles using interests and behaviors, generates interaction trajectories, annotates each interaction record using the DPS vector, and estimates exposure as \(E(D \mid A)\). By comparing exposure estimates across profiles, the framework examines whether dark pattern exposure is associated with profile differences. This approach does not claim to reveal the internal personalization logic of mobile applications, nor does it establish direct causal intent. Instead, it provides an empirical method for observing and comparing the interface exposure experienced by controlled profiles.

The next chapter builds on this related work by defining the conceptual framework of profile-aware dark pattern auditing. It formalizes user profiles, interaction trajectories, DPS annotation, and profile-conditioned exposure estimation, thereby translating the research gap identified here into an operational audit design.
